% !TeX root = ../../../main.tex
\section{\'Etale Torsors, Covers, and Quotients}
Fix a scheme $X$, and let $G$ be a smooth affine $X$-group scheme. 
Denote by $X_{\text{Ét}}$ and $X_{\text{ét}}$ the big and small étale sites on $X$, respectively. 
The functor of points of $G$, which we denote by $G(\cdot) = \mathrm{Hom}_X(-, G)$, defines a group sheaf on 
$X_{\text{ét}}$.

\begin{definition} A \textit{principal $G$-bundle} (or \textit{principal homogeneous space}) for $G$ is a 
scheme $f : Y \to X$ equipped with a left $G$-action $\rho: G \times_X Y \to Y$ satisfying:
\begin{enumerate}
    \item The morphism $G \times_X Y \to Y \times_X Y$ sending $(g, y) \mapsto (gy, y)$ is an isomorphism of $X$-schemes.
    \item There exists an étale covering $\{U_i \to X\}$ such that $Y_{U_i} \cong G_{U_i}$ as $G$-schemes, 
    where $G$ acts on itself by left translation.
\end{enumerate}
\end{definition}

\begin{rem*}
    Since $G \to X$ is smooth and $Y$ is étale-locally isomorphic to $G$, 
    the morphism $f: Y \to X$ is smooth and affine. Note that $Y$ is generally not étale over $X$ unless $G$ is an 
    étale group scheme (e.g., a finite constant group).
    Similarly, if $G$ is finite then $Y \to X$ is finite.
\end{rem*}

Similarly to the case of $G$-torsors of a topos, we define a morphism of principal $G$-bundles to be a morphism of $X$-schemes commuting with the
$G$-action. We now have:

\begin{theorem}\label{theorem:torsors_are_realized_as_spaces}
    The morphism sending $Y \mapsto \mathrm{Hom}_X(-, Y)$ is an equivalence of categories from the category of principal $G$-bundles 
    to the category of $G$-torsors on $X_{\text{\'Et}}$
    Similarly, the morphism sending $Y$ to $\mathrm{Hom}_X(-, Y)$ to the category of $G$-torsors on $X_{\text{ét}}$ is an equivalence.
\end{theorem}

%% [REVIEW: unlabeled & never cited -- the big/small-site comparison is used implicitly whenever torsors are realized as schemes; consider labeling it and citing it (e.g. line ~381), or merging into the theorem above]
\begin{cor}
    There is a natural equivalence $\mathbf{Tors}(X_{\mathrm{\text{\'et}}}, G) \cong \mathbf{Tors}(X_{\text{\'Et}}, G)$ 
    inducing a bijection of pointed sets $\mathrm{Tors}(X_{\text{\'et}}, G) \xrightarrow{\cong} \mathrm{Tors}(X_{\text{\'Et}}, G)$ 
    which is an isomorphism of abelian groups if $G$ is abelian.
    And every $G$-torsor, in each of the above sites can be realized as a scheme, which is a prinicipal $G$-bundle.
\end{cor}

%%% %%% %%% %%% %%% %%% %%% %%%
%       Torsors and Cohomology
%%% %%% %%% %%% %%% %%% %%% %%%

% \subsection*{Torsors and Cohomology}
% We qoute without proof:
% \begin{theorem}
% There is a natural bijection of pointed sets     
% $\text{Tors}(\mathcal{G|_T}) \to \check{H}^1(T, \mathcal{G})$. Moreover, if $\mathcal{G}$ is an abelian group sheaf, it's an isomorphism of abelian groups.
% \end{theorem}

% The idea is that for every $[\mathcal{P}] \in \text{Tors}(T, \mathcal{G})$ 
% We 
% \begin{enumerate}
%     \item Choose a covering $\{U_i \to T\}$ such that $\mathcal{P}(U_i) \neq \emptyset$ for all $i$
%     \item choose sections $\alpha_i \in \mathcal{F}(U_i)$
% \end{enumerate}

% Then we note that for all $(i,j)$ the elements $\alpha_i |_{U_i \times_X U_j}$ and $\alpha_j |_{U_i \times_X U_j}$ differ by a unique element of $\mathcal{G}(U_{ij})$
% there exists a unique $s_{ij} \in \mathcal{G}(U_{ij})$ such that $\alpha_i |_{U_i \times_X U_j} = s_{ij}(\alpha_j |_{U_i \times_X U_j})$. 
% One can then easily see that $(s_{ij})$ defines an element of  $\check{H}^1(T, \mathcal{G})$
% which is independent of the choice of repreantiative $\cP$, choice of covering $\{U_i \to T\}$ and choice of sections.

We will also use the standard identification between torsors and first
\'etale cohomology. For an abelian sheaf of groups $G$ on
$X_{\mathrm{\acute et}}$, there is a natural group isomorphism
\[
\Tors(X_{\mathrm{\acute et}},G)\xrightarrow{\ \cong\ }
H^1_{\mathrm{\acute et}}(X,G),
\]
where the group operation on torsor classes is induced by contracted
product \cite[III, \S4, Proposition~4.6 and Remark~4.8(b)]{milneEC}. This identification allows us to describe
cyclic torsors by their cohomology classes in Kummer, Artin--Schreier, and
Artin--Schreier--Witt theory.



%%% %%% %%% %%% %%% %%% %%% %%%
%       Constant Finite Group Torsors
%%% %%% %%% %%% %%% %%% %%% %%%

\subsection*{Constant Finite Group Torsors}

Let $G$ be a finite group, within this section, we denote the associated constant group scheme over $X$ by $\underline{G}$ to maintain a clear distinction between the group as a set and the group as a scheme. In subsequent sections, we shall follow the standard practice of identifying $G$ with its constant group scheme and omit the underline for brevity.

$\underline{G}$ is given by $\coprod_{g \in G} X$ with the action shuffeling the $X$'s according to multiplication.

By following the definitions, one sees that if $X$ be a connected scheme and $f : Y \to X$ is a finite Galois cover with Galois group $G$. 
Then, $f : Y \to X$ is a principal $\underline{G}$-bundle.
(Recall that a \textit{finite Galois cover} is a finite étale surjection $Y \to X$ with $Y$ connected and such that 
$G = \text{Aut}(Y/X)$ acts transitively on the geometric points of $Y$ lying over any geometric point of $X$).

On the otherhand if $f: Y \to X$ is a principal $\underline{G}$-bundle with $Y$ connected, then $Y$ is a finite Galois cover with
automorphism group $G$.

However, not all $\underline{G}$-torsors are connected. If $H\subset G$ is a proper subgroup then any connected finite \'etale cover $f: Y \to X$ with 
Galois group $H$ gives rise to a non-connected $\underline{G}$-torsor by looking at the induced $\underline{G}$-torsor $\varphi_*(Y)$ under the inclusion $\varphi: \underline{H} \to \underline{G}$.
On the otherhand, if we fix a geometric point $\overline{x} \to X$, then to give an homomorphism
$\rho \in \mathrm{Hom}_{\mathrm{cont}}(\pi_1^{\mathrm{et}}(X, x), G)$ is equivilant to give a connected pointed Galois cover
$(Y, \overline{y}) \to (X, \overline{x})$ with Galois group $H=\rho(\pi_1^{\mathrm{et}}(X, x)) \subset G$. 
Thus, pushing forward to $\underline{G}$ we get a prinicipal $\underline{G}$-bundle. The choice of a different geometric point $\overline{x'} \to X$
differ the homomorphism by an inner automorphism, thus we have:

\begin{theorem}\label{theorem:equivalence_fundamental_covers}
    Let $X$ be a connected scheme and $\overline{x}$ a geometric point of $X$. Suppose in addition that $G$ is a finite abstract group. 
    Define a map
    \begin{equation}
        \mathrm{Hom}_{\mathrm{cont}}(\pi_1^{\mathrm{et}}(X, \overline{x}), G) / \mathrm{Inn}(G) \to \mathrm{Tors}(X_{\text{\'Et}}, \underline{G}) 
    \end{equation}
by sending a homomorphism $\rho : \pi_1^{\mathrm{et}}(X, \overline{x}) \to G$ 
to the principal $\underline{G}$-bundle $\varphi_*(Y)$ where $Y$ is the principal $\underline{\rho(\pi_1^{\mathrm{et}}(X, \overline{x}))}$-bundle 
obtained above and $\varphi$ is the inclusion $\rho(\pi_1^{\mathrm{et}}(X, \overline{x})) \hookrightarrow G$. Then, the map 
 is a bijection of pointed sets where the trivial homomorphisms
  (which is the only element of its $\mathrm{Inn}(G)$-orbit) is the distinguished element of the left hand side.
\end{theorem}

If $G$ is abelian then $\mathrm{Inn}(G)$ is trivial, and we obtain:
\begin{cor}\label{cor:abelian_equivilance_fundamental_torsors}
    Let $G$ be a finite abelian group, $X$ a connected scheme, and $\overline{x}$ a geometric point of $X$. Then, the map from \Cref{theorem:equivalence_fundamental_covers}
     induces an isomorphism of abelian groups
    \begin{equation}
        \mathrm{Hom}_{\mathrm{cont.}}(\pi_1^{\text{ét}}(X, \overline{x}), G) \xrightarrow{\cong} \Tors(X_{\text{\'Et}}, \underline{G})
    \end{equation}
\end{cor}

For $G=\Lambda^\times$, the character in
\Cref{cor:abelian_equivilance_fundamental_torsors} is also the monodromy
character of the associated rank-one $\Lambda$-local system. In particular,
the torsor and the local system have the same restrictions to inertia and to
the upper ramification groups. Thus, one has ramification bounded by a given
modulus, or is tamely ramified along a given divisor, if and only if the other
does.

\subsection*{Quotients and Decomposition of Torsors}

Let $\mathcal{H}$ be a sheaf of groups acting on a sheaf $\mathcal{P}$ on a
site $\sC$. The \textbf{quotient sheaf} $\mathcal{P}/\mathcal{H}$ is the
coequalizer in $\operatorname{Sh}(\sC)$
\[
\mathcal{H}\times\mathcal{P}
\rightrightarrows
\mathcal{P}\longrightarrow\mathcal{P}/\mathcal{H},
\]
where the parallel arrows are the action and the projection.
Equivalently, it is the sheaf associated with the presheaf of orbits
$U\mapsto\mathcal{H}(U)\backslash\mathcal{P}(U)$. We use the notation
$\mathcal{P}/\mathcal{H}$ even though our torsor actions are written on the
left. In particular, the quotient always exists as a sheaf. It is
characterized by the following universal property:
for every sheaf $\mathcal{F}$, regarded as an $\mathcal{H}$-sheaf with the
trivial action, composition with
$\mathcal{P}\to\mathcal{P}/\mathcal{H}$ identifies morphisms
$\mathcal{P}/\mathcal{H}\to\mathcal{F}$ with the
$\mathcal{H}$-equivariant morphisms $\mathcal{P}\to\mathcal{F}$, equivalently
with morphisms $f$ satisfying $f(hp)=f(p)$ on local sections. A quotient
sheaf need not in general be representable by a scheme; representability in
the finite-torsor situation used here is part of the next proposition.

Let $G$ be a finite abelian group, let $H\subseteq G$ be a subgroup, and let
$\pi:G\to G/H$ be the quotient homomorphism.

\begin{prop}\label{prop:quotient_torsors_finite}\label{prop:quotient_torsors}
Let $\mathcal{P}\to X$ be a $G$-torsor in $X_{\text{ét}}$. The quotient sheaf
$\mathcal{P}/H$ is represented by a finite étale $X$-scheme, and there is a
canonical isomorphism of $G/H$-torsors
\[
\mathcal{P}/H \xrightarrow{\sim}
\pi_*\mathcal{P}=(G/H)\times^G\mathcal{P}.
\]
Consequently, $\mathcal{P}\to X$ admits a canonical factorization
\[
\mathcal{P}\xrightarrow{\phi}\mathcal{P}/H\xrightarrow{\psi}X,
\]
where $\phi$ is an $H$-torsor and $\psi$ is a $G/H$-torsor. This
factorization is unique up to a unique isomorphism commuting with the
projections.
\end{prop}

\begin{proof}
The canonical morphism
\[
\mathcal{P}\longrightarrow\pi_*\mathcal{P},
\qquad p\longmapsto[1,p],
\]
is $H$-invariant because $\pi(H)=1$. The universal property of the quotient
sheaf therefore gives a canonical morphism
\[
\mathcal{P}/H\longrightarrow\pi_*\mathcal{P}.
\]
Choose an étale cover $U\to X$ that trivializes $\mathcal{P}$. Over $U$ this
morphism identifies with the identity of $(G/H)\times U$, since
\[
(G\times U)/H\cong(G/H)\times U
\qquad\text{and}\qquad
\pi_*\mathcal{P}|_U\cong(G/H)\times U.
\]
It is therefore an isomorphism of sheaves. The contracted product
$\pi_*\mathcal{P}$ is a $G/H$-torsor and, by
\Cref{theorem:torsors_are_realized_as_spaces}, is represented by a finite
étale $X$-scheme. Hence the same is true of $\mathcal{P}/H$.

Under the same local trivialization, the morphism
$\mathcal{P}\to\mathcal{P}/H$ is the product of the quotient map with the
identity,
\[
G\times U\longrightarrow(G/H)\times U,
\]
and is thus an $H$-torsor. Finally, suppose
$\mathcal{P}\to\mathcal{Q}\to X$ is another such factorization. Since an
$H$-torsor is locally split, every $H$-invariant morphism out of
$\mathcal{P}$ descends uniquely through $\mathcal{Q}$. Thus $\mathcal{Q}$
satisfies the same universal property as $\mathcal{P}/H$, which gives the
unique isomorphism compatible with the projections.
\end{proof}

